\subsection{BÀI TẬP TRẮC NGHIỆM}
\ind{PHẦN I.} \inden{Câu trắc nghiệm nhiều phương án lựa chọn. Học sinh trả lời từ câu 1 đến câu 16. Mỗi câu hỏi học sinh chỉ chọn một phương án.}\\
\setcounter{ex}{0}
\Opensolutionfile{ans}[ans/1D2-B1-1]
\begin{ex}%[1D3Y2-2]
	Cho dãy số $({{u}_{n}} ),$ biết ${{u}_{n}}={{(-1 )}^{n}}\cdot 2n$. Mệnh đề nào sau đây {\bf sai}?
	\choice
	{${{u}_{3}}=-6$}
	{${{u}_{2}}=4$}
	{ \True ${{u}_{4}}=-8$}
	{${{u}_{1}}=-2$}
	\loigiai {
		Thay trực tiếp hoặc dùng chức năng CALC:\\
		${{u}_{1}}=-2\cdot 1=-2;\,\,{{u}_{2}}={{(-1 )}^2}\cdot 2\cdot 2=4,\,\,{{u}_{3}}={{(-1 )}^3}\cdot 2\cdot 3=-6;\,\,{{u}_{4}}={{(-1 )}^4}\cdot 2\cdot 4=8$.\\
		\textbf{Nhận xét:} Dễ thấy ${{u}_{n}}>0$ khi $n$ chẵn và ngược lại nên đáp án $u_4=-8$ sai.}
\end{ex}

\begin{ex}%[1D3Y2]
	Cho dãy số $(u_n)$ với $u_n=2n+1$. Tìm $u_5$.
	\choice
	{ \True $11$}
	{ $2$}
	{ $1$} 
	{ $3$}
	\loigiai{
		Với $n=5\Rightarrow u_5=11.$}
\end{ex}


\begin{ex}%[1D3Y2-2]
	Cho dãy số $({{u}_{n}} )$, biết ${{u}_{n}}=\dfrac{n}{{{3}^{n}}-1}$. Ba số hạng đầu tiên của dãy số đó lần lượt là những số nào dưới đây?
	\choice
	{$\dfrac{1}{2};\dfrac{1}{4};\dfrac{1}{16}$}
	{$\dfrac{1}{2};\dfrac{2}{3};\dfrac{3}{4}$}
	{ \True $\dfrac{1}{2};\dfrac{1}{4};\dfrac{3}{26}$}
	{$\dfrac{1}{2};\dfrac{1}{4};\dfrac{1}{8}$}
	\loigiai {
		Dùng MTCT chức năng CALC: ta có
		${{u}_{1}}=\dfrac{1}{2};\,\,{{u}_{2}}=\dfrac{2}{{{3}^2}-1}=\dfrac{2}{8}=\dfrac{1}{4};\,\,{{u}_{3}}=\dfrac{3}{{{3}^3}-1}=\dfrac{3}{26}$.}
\end{ex}


\begin{ex}%[1D3B2]
	Cho dãy số $\left(u_n\right)$ xác định bởi $\heva{& u_1=3 \\& u_{n+1}=\sqrt{1+u_n^2} \,\text{ với } n\ge 1}$. Tìm số hạng thứ hai của dãy số $\left(u_n\right)$.
	\choice
	{$u_2=2$}
	{\True $u_2=\sqrt{10}$}
	{$u_2=10$}
	{$u_2=\sqrt{2}$}
	\loigiai{
		Ta có $\heva{& u_1=3 \\& u_2=u_{1+1}=\sqrt{1+u_1^2}=\sqrt{10}.}$
	}
\end{ex}

\begin{ex}%[1D3Y2]%
	Dãy $(u_n)$ gồm có 5 phần tử cho bởi $\begin{cases} u_1 = 1 \\ u_{n+1} = u_n + 2, \ \forall n \ge 1. \end{cases}$ Phần tử thứ 5 của dãy bằng
	\choice
	{$7$}
	{$5$}
	{\True $9$}
	{$3$}
	\loigiai{
		{\bf Cách 1:} (nếu chưa học tới cấp số cộng) Đây là dãy số cho bởi công thức truy hồi, ta cứ thay vào công thức thì được 
		$$u_1 = 1;\ u_2=3; \ u_3 = 5; \ u_4 = 7; \ u_5 = 9.$$
		{\bf Cách 2:} (nếu đã học tới cấp số cộng) Đây là cấp số cộng có $u_1=1$ và công sai $d=2$ nên ta có $$u_5=u_1 + 4d = 1+4.2=9.$$
	}
\end{ex}

\begin{ex}%[1D3Y2-2]
	Cho dãy số $(u_n):\heva{&u_1=u_2=1\\&u_n=u_{n-1}+u_{n-2}, \forall n\geq 3.}$ Tìm số hạng thứ $7$ của dãy.
	\choice
	{\True $u_7=13$}
	{$u_7=21$}
	{$u_7=17$}
	{$u_7=7$}
	\loigiai{
		Ta có $u_3=u_1+u_2=2, u_4=u_2+u_3=3, u_5=u_3+u_4=5, u_6=u_4+u_5=8, u_7=u_5+u_6=13$.
	}
\end{ex}

\begin{ex}%[1D3B2]
	Cho dãy số $\left(u_n\right)$ với $u_n=\dfrac{3n-4}{2n+5}$. Số $\dfrac{14}{17}$ là số hạng thứ bao nhiêu của dãy? 	
	\choice
	{$4$}
	{$7$}
	{$5$}
	{\True $6$}
	\loigiai{Ta có $\dfrac{3n-4}{2n+5}=\dfrac{14}{17}\Leftrightarrow 23n=138 \Leftrightarrow n=6$.\\
		Vậy $\dfrac{14}{17}$ là số hạng thứ $6$ của dãy.
	}
\end{ex}

\begin{ex}%[1D3B2-2]
	Cho dãy số $({{u}_{n}} ),$ biết ${{u}_{n}}=\dfrac{n+1}{2n+1}$. Số $\dfrac{8}{15}$ là số hạng thứ mấy của dãy số?
	\choice
	{ \True $7$}
	{$6$}
	{$5$}
	{$8$}
	\loigiai {
		Ta cần tìm $n$ sao cho ${{u}_{n}}=\dfrac{n+1}{2n+1}=\dfrac{8}{15}\Leftrightarrow 15n+15=16n+8\Leftrightarrow n=7$.\\
		\textbf{Nhận xét:} Có thể dùng chức năng CALC để kiểm tra nhanh.}
\end{ex}

\begin{ex}%[Đề thi HK1 THPT Hoài Đức A, Hà Nội, 2018]%[1D3Y2]
	Cho các dãy số sau, dãy số nào là dãy tăng?
	\choice{$1$; $-\dfrac{1}{2}$; $\dfrac{1}{3}$; $-\dfrac{1}{4}$}
	{$1$; $\dfrac{1}{2}$; $\dfrac{1}{4}$; $\dfrac{1}{6}$}
	{$-1$; $3$; $5$; $3$}
	{\True $2$; $4$; $6$; $8$}
	\loigiai{}
\end{ex}

\begin{ex}%[Đề thi học kỳ 1, Sở Bà Rịa Vũng Tàu, 2018-2019]%[Đỗ Viết Lân]%[1D3B2-3]
	Dãy số nào có công thức số hạng tổng quát dưới đây là dãy số tăng?
	\choice
	{$u_n=\left(\dfrac{1}{2}\right)^n$}
	{$u_n=( - 3)^n$}
	{$u_n=2020 - 3n$}
	{\True $u_n=2018 + 2n$}
	\loigiai{
		Xét từng dãy số
		\begin{itemize}
			\item $u_n=\left(\dfrac{1}{2}\right)^n$, ta có $u_1 = \dfrac{1}{2}$ và $u_2 = \dfrac{1}{4}$ nên dãy này không tăng.
			\item $u_n=( - 3)^n$, ta có $u_2=9$ và $u_3 = - 27$ nên dãy này không tăng.
			\item $u_n=2020 - 3n$, ta có $u_1 = 2017$ và $u_2 = 2014$ nên dãy số không tăng.
			\item $u_n=2018 + 2n$, ta có $u_{n+1} = 2020+2n$ suy ra $u_{n+1}-u_n = 2 > 0$. Do đó dãy này là dãy số tăng.
		\end{itemize}
	}
\end{ex}

\begin{ex}%[Thi học kỳ 1, Dĩ An, Bình Dương, 2018]%[1D3B2]
	Trong các dãy số $\left(u_n\right)$ sau đây, hãy chọn dãy số giảm.  
	\choice
	{\True $u_n=\sqrt{n}-\sqrt{n-1}$}
	{$u_n=\sin n$}
	{$u_n=\dfrac{n^2+1}{n}$}
	{$(-1)^n\left(2^n+1\right)$}
	\loigiai{
		Xét dãy số $\left(u_n\right)$ với $u_n=\sqrt{n}-\sqrt{n-1}$, vì $\sqrt{n}>\sqrt{n-1}$ với $n\in\mathbb{N}^*$ nên $u_n>0,\forall n\in\mathbb{N}^*$, ta có:\\
		$\dfrac{u_{n+1}}{u_n}=\dfrac{\sqrt{n+1}-\sqrt{n}}{\sqrt{n}-\sqrt{n-1}}=\dfrac{(n+1-n)\left(\sqrt{n}+\sqrt{n-1}\right)}{\left[n-(n-1)\right]\left(\sqrt{n+1}+\sqrt{n}\right)}=\dfrac{\sqrt{n}+\sqrt{n-1}}{\sqrt{n+1}+\sqrt{n}}$.\\
		Ta thấy $\sqrt{n+1}>\sqrt{n-1}\Rightarrow \sqrt{n+1}+\sqrt{n}>\sqrt{n}+\sqrt{n-1}$, suy ra $\dfrac{u_{n+1}}{u_n}<1$ hay $u_{n+1}<u_n$.\\
		Vậy dãy số $\left(u_n\right)$ với $u_n=\sqrt{n}-\sqrt{n-1}$ là dãy số giảm.
	}
\end{ex}

\begin{ex}%[1D3B2]
	Trong các dãy số $\left(u_n\right)$ sau, dãy số nào là dãy số bị chặn?
	\choice
	{\True $u_n=\dfrac{n}{n+1}$}
	{$u_n=\sqrt{n^2+1}$}
	{$u_n=2^n+1$}
	{$u_n=n+\dfrac{1}{n}$}
	\loigiai{
		Xét $u_n=\dfrac{n}{n+1}=1-\dfrac{1}{n+1}$ $\Rightarrow 0<u_n<1$, $\forall n \in \mathbb{N}^*$.\\
		Vậy $u_n=\dfrac{n}{n+1}$ là dãy số bị chặn.
	}
\end{ex}

\begin{ex}%[Thi HK1, THPT Lương Thế Vinh Hà Nội, 2019]%[Tuan Nguyen, dự án EX11]%[1D3B2-1]
	Cho dãy số $(u_n)$ biết $u_n=2^n$. Mệnh đề nào sau đây là mệnh đề đúng?
	\choice
	{$u_{n+2}=2^2$}
	{$u_{n+2}=2^n+2$}
	{$u_{n+2}=2\cdot 2^{n}$}
	{\True $u_{n+2}=4\cdot 2^{n}$}
	\loigiai{
		Ta có $u_n=2^n\Rightarrow u_{n+2}=2^{n+2}=4\cdot 2^n$.
	}
\end{ex}

\begin{ex}%[SỞ GD \& ĐT BÌNH ĐỊNH]%[1D3B2]
	Cho dãy $ (u_n) $ là: $ 0;\; \dfrac{1}{2}; \; \dfrac{2}{3}; \; \dfrac{3}{4}; \; \dfrac{4}{5};\ldots $. Số hạng tổng quát $ (u_n) $  là
	\choice
	{$ \dfrac{n+1}{n} $}
	{\True $ \dfrac{n-1}{n} $}
	{ $ \dfrac{n}{n+1} $}
	{$ \dfrac{n^2-n}{n+1} $}
	\loigiai{
		
	}
\end{ex}

\begin{ex}%[1D3B2-1]
	Cho dãy số có các số hạng đầu là: $-2;0;2;4;6;\ldots$. Số hạng tổng quát của dãy số này là công thức nào dưới đây?
	\choice
	{${{u}_{n}}=-2(n+1 )$}
	{${{u}_{n}}=-2n$}
	{\True ${{u}_{n}}=2n-4$}
	{${{u}_{n}}=n-2$}
	\loigiai {
		Kiểm tra ${{u}_{1}}=-2$ ta loại các đáp án ${{u}_{n}}=n-2$, ${{u}_{n}}=-2(n+1 )$. Ta kiểm tra ${{u}_{2}}=0$\\
		Xét đáp án ${{u}_{n}}=-2n$
		${{u}_{n}}=n-2$: ${{u}_{n}}=2n\Rightarrow  {{u}_{2}}=4 \neq 0\Rightarrow {{u}_{n}}=n-2$ loại.\\
		Xét đáp án ${{u}_{n}}=2n-4$: ${{u}_{n}}=2n-4=2\cdot 2-4=0$ nhận.\\
		Nhận xét: Dãy $2;\,4;\,6$ có công thức là $2n\,\,(n\in {{\mathbb{N}}^{*}} )$ nên dãy $-2;0;2;4;6;\ldots $có được bằng cách “tịnh tiến” $2n$ sang trái 4 đơn vị, tức là $2n-4$}
\end{ex}

\begin{ex}%[1D3B2-1]
	Cho dãy số $({{u}_{n}} ),$ được xác định $\heva{
		& {{u}_{1}}=2 \\
		& {{u}_{n+1}}=2u_n.}$ Số hạng tổng quát ${{u}_{n}}$ của dãy số là số hạng nào dưới đây?
	\choice
	{ \True ${{u}_{n}}={{2}^{n}}$}
	{${{u}_{n}}={{2}^{n+1}}$}
	{${{u}_{n}}=2$}
	{${{u}_{n}}={{n}^{n-1}}$}
	\loigiai {
		Từ công thức $\heva{
			& {{u}_{1}}=2 \\
			& {{u}_{n+1}}=2{{u}_{n}} \\
		}\Rightarrow \heva{
			& {{u}_{1}}=2 \\
			& {{u}_{2}}=2{{u}_{1}}=2\cdot 2=4 \\
			& {{u}_{3}}=2{{u}_{2}}=2\cdot 4=8 \\
		}.$\\
		Xét đáp án ${{u}_{n}}={{n}^{n-1}}$ với $n=1\Rightarrow {{u}_{1}}={{1}^{1-1}}={{1}^{0}}=1$ loại.\\
		Xét đáp án ${{u}_{n}}={{2}^{n}}$, ta thấy đều thỏa mãn. \\
		Xét đáp án ${{u}_{n}}={{2}^{n+1}}$ với $n=1\Rightarrow{{u}_{1}}={{2}^{1+1}}={{2}^2}=4$ loại.\\
		Dễ thấy đáp án ${{u}_{n}}=2$ không thỏa mãn}
\end{ex}



\Closesolutionfile{ans}

\ind{PHẦN II.} \inden{Câu trắc nghiệm đúng sai. Học sinh trả lời từ câu 17 đến câu 20. Trong mỗi ý a), b), c), d) ở mỗi câu, học sinh chọn đúng hoặc sai.}\\
\Opensolutionfile{ans}[ans/1D2-B1-2]

\begin{ex}%[1D2B1-4]
	Cho dãy số $\left(u_n\right)$ với $u_n=\dfrac{2 n+1}{n+2}$. Xét tính đúng sai của các khẳng định sau:
	\choiceTF
	{\True $u_3+u_5=\dfrac{104}{35}$}
	{\True $\left(u_n\right)$ có đúng một số hạng là số nguyên}
	{$\left(u_n\right)$ là dãy số giảm}
	{\True $\left(u_n\right)$ là dãy số bị chặn}
	\loigiai{
			\begin{itemchoice}
			\itemch Ta có $u_3=\dfrac{7}{5}$, $u_5=\dfrac{11}{7}$. Suy ra $u_3+u_5=\dfrac{7}{5}+\dfrac{11}{7}=\dfrac{104}{35}$.
			\itemch Xét $u_n=\dfrac{2 n+1}{n+2}=2-\dfrac{3}{n+2}$. Để $u_n$ là số hạng nguyên thì $3$ phải chia hết cho $n+2$, suy ra $n=1$. Vậy $\left(u_n\right)$ có đúng một số hạng là số nguyên là $u_1=1$.
			\itemch Xét $u_{n+1}-u_n=2-\dfrac{3}{n+3}-\left( 2-\dfrac{3}{n+2}\right)=\dfrac{3}{n+2}-\dfrac{3}{n+3}>0, \forall n \in \mathbb{N}^*$. Suy ra $(u_n)$ là dãy tăng.
			\itemch Ta có $u_n=\dfrac{2 n+1}{n+2}>0, \forall n \in \mathbb{N}^*$, suy ra $\left(u_n\right)$ bị chặn dưới.\\
			Ta lại có $u_n=\dfrac{2 n+1}{n+2}=\dfrac{2(n+2)-3}{n+2}=2-\dfrac{3}{n+2}<2, \forall n \in \mathbb{N}^*$, suy ra $\left(u_n\right)$ bị chặn trên.\\
			Dãy số $\left(u_n\right)$ vừa bị chặn trên vừa bị chặn dưới nên dãy số $\left(u_n\right)$ bị chặn.
		\end{itemchoice}
	}
\end{ex}

\begin{ex}%[1D3K2-4]
	Cho dãy số $(u_n ),$ với $u_n=2024^n$. Xét tính đúng sai của các khẳng định sau:
	\choiceTF
	{\True $\dfrac{u_2}{u_1}=2024$}
	{Số hạng thứ $n+1$ là $u_{n+1}=2024^n+1$}
	{\True Dãy số $(u_n)$ là một dãy số tăng}
	{Dãy số $(u_n)$ là dãy số bị chặn}
	\loigiai {
		\begin{itemchoice}
			\itemch Ta có $\dfrac{u_2}{u_1}=\dfrac{2024^2}{2024^1}=2024$.
			\itemch Số hạng thứ $n+1$ là $u_{n+1}=2024^{n+1}$.
			\itemch Ta có $\dfrac{u_{n+1}}{u_n}=\dfrac{2024^{n+1}}{2024^n}=2024>0$ nên $(u_n)$ là dãy số tăng
			\itemch Dãy $(u_n)$ không có chặn trên nên $(u_n)$ không bị chặn.
		\end{itemchoice}
	}
\end{ex}

\begin{ex}%[1D2K1-2]
	Cho dãy số $\left(u_n\right)$, biết $u_n=\sin \left[(2 n-1) \dfrac{\pi}{4}\right]$. Xét tính đúng sai của các khẳng định sau:
	\choiceTF
	{Số hạng đầu tiên của dãy là $u_1=1$}
	{Dãy $(u_n)$ là dãy số tăng}
	{\True $u_{n+4}=u_n$, với mọi $n \geq 1$}
	{\True Tổng 12 số hạng đầu của dãy số bằng 0}
	\loigiai{
		\begin{itemchoice}
			\itemch $u_1=\sin \left[(2\cdot 1-1)\dfrac{\pi}{4}\right]=\sin \dfrac{\pi}{4}=\dfrac{\sqrt{2}}{2}$.
			\itemch Do $u_1=\dfrac{\sqrt{2}}{2}$, $u_3=-\dfrac{\sqrt{2}}{2}$ nên $(u_n)$ không là dãy số tăng.
			\itemch Với mọi $n\ge 1$, ta có 
				\begin{eqnarray*}
				u_{n+4}
				&=&\sin \left[\left[2(n+4)-1\right]\dfrac{\pi}{4}\right]=\sin\left[(2n-1)\dfrac{\pi}{4}+2\pi\right]\\
				&=&  \sin\left[(2n-1)\dfrac{\pi}{4}\right]=u_n.
			\end{eqnarray*}
			\itemch Gọi $S_{12}$ là tổng $12$ số hạng đầu của dãy số.
			\begin{eqnarray*}
				S_{12}&=&\sin\dfrac{\pi}{4}+\sin\dfrac{3\pi}{4}+\sin\dfrac{5\pi}{4}+\sin\dfrac{7\pi}{4}+\cdots+\sin\dfrac{23\pi}{4}\\
				&=&\left(\sin\dfrac{\pi}{4}+\sin\dfrac{23\pi}{4}\right)+\left(\sin\dfrac{3\pi}{4}+\sin\dfrac{21\pi}{4}\right)+\cdots+\left(\sin\dfrac{11\pi}{4}+\sin\dfrac{13\pi}{4}\right)\\
				&=&\left[\sin\dfrac{\pi}{4}+\sin\left(6\pi-\dfrac{\pi}{4}\right)\right]+\left[\sin\dfrac{3\pi}{4}+\sin\left(6\pi-\dfrac{3\pi}{4}\right)\right]+\cdots+\left[\sin\dfrac{11\pi}{4}+\sin\left(6\pi-\dfrac{11\pi}{4}\right)\right]\\
				&=&\left(\sin\dfrac{\pi}{4}-\sin\dfrac{\pi}{4}\right)+\left(\sin\dfrac{3\pi}{4}-\sin\dfrac{3\pi}{4}\right)+\cdots+\left(\sin\dfrac{11\pi}{4}-\sin\dfrac{11\pi}{4}\right)\\
				&=&0.
			\end{eqnarray*}
		\end{itemchoice}
		}
\end{ex}

\begin{ex}%[1D2B1-5]
		Gọi $u_n$ là tổng diện tích của các hình vuông có ở hàng thứ $n$ trong hình bên (mỗi ô vuông nhỏ là 1 đơn vị diện tích). Xét tính đúng sai của các khẳng định sau:
		\begin{center}
			\begin{tikzpicture}[line join=round, line cap = round, >=stealth, scale=1,font=\footnotesize,transform shape]
			\foreach \i in {1,2,3,4}
			{
				\draw (0,4.5 - \i) node{$\text{Hàng thứ }\i$};
				\pgfmathsetmacro\x{.1*\i}
				\foreach \j in {1,...,\i}
				{
					\begin{scope}[shift={(1 + \j,4.5 - \i - \x)}]
						\draw[fill = cyan!70,scale=.2] (0,0) rectangle (\i,\i);
						\draw[scale=.2] (0,0) grid (\i,\i);
					\end{scope}
				}
			}
		\end{tikzpicture}
		\end{center}
		\choiceTF
		{\True $u_1=1$}
		{$u_2=4$}
		{\True $u_1+u_2+u_3+u_4=100$}
		{$u_n = n^2$ với $n\in\mathbb{N}^*$}	
	\loigiai{
		Ta có 
		\begin{itemize}
			\item Hàng $1$ có $1$ hình vuông cạnh $1$ đơn vị nên $u_1 = 1\times 1 = 1^3$.
			\item Hàng $2$ có $2$ hình vuông cạnh $2$ đơn vị nên $u_2 = 2\times 2^2 = 2^3$.
			\item Hàng $3$ có $3$ hình vuông cạnh $3$ đơn vị nên $u_3 = 3\times 3^2 = 3^3$.
			\item Hàng $4$ có $4$ hình vuông cạnh $4$ đơn vị nên $u_4 = 4\times 4^2 = 4^3$.
		\end{itemize}
	nên
		\begin{itemchoice}
			\itemch $u_1 = 1$.
			\itemch $u_2 = 8$.
			\itemch $u_1 = 1; u_2 = 8; u_3 = 27; u_4 = 64$, suy ra $u_1+u_2+u_3+u_4=100$.
			\itemch Công thức tổng quát của dãy số $(u_n)$ là $u_n = n^3$ với $n\in\mathbb{N}^*$.
		\end{itemchoice}
	}
\end{ex}
\Closesolutionfile{ans}

\ind{PHẦN III.} \inden{Câu trắc nghiệm trả lời ngắn. Học sinh trả lời từ câu 21 đến câu 25 vào ô kết quả.}\\
\Opensolutionfile{ans}[ans/1D2-B1-3]

\begin{ex}
	Cho dãy số $(u_n)$ có $u_1=2$, $u_2=3$ và $u_{n+1}=2u_n+u_{n-1}$ với mọi $n\ge 2$, $n\in\mathbb{N}$. Tìm số hạng thứ tư của dãy số đó.\\
	\shortans[oly]{19}
	\loigiai{
		Từ $u_3=2\cdot 3+2=8\Rightarrow u_4=2\cdot 8+3=19$.}
\end{ex}



\begin{ex}
	Nếu tỉ lệ lạm phát là $3,5 \%$ mỗi năm và giá trung bình của một căn hộ chung cư mới tại thời điểm hiện tại là $2,5$ tỉ đồng thì giá trung bình của một căn hộ chung cư mới sau $n$ năm nữa được cho bởi công thức
	$$
	A_{n}=2,5 \cdot(1,035)^{n} \text { (tỉ đồng). }
	$$
	Tìm giá trung bình (tỉ đồng)của một căn hộ chung cư mới sau $5$ năm nữa (kết quả làm tròn đến hàng phần trăm).\\
		\shortans[oly]{$2,97$}
	\loigiai{
		Giá trung bình của một căn hộ chung cư mới sau $5$ năm là
		$$
		A_{5}=2,5 \cdot(1,035)^{5}=2,9692 \text { (tỉ đồng). }
		$$
	}
\end{ex}

\begin{ex}
	Giá của một chiếc máy photocopy lúc mới mua là 50 triệu đồng. Biết rằng giá trị của nó sau mỗi năm sử dụng chỉ còn $75 \%$ giá trị trong năm liền trước đó. Tính giá trị còn lại (triệu đồng) của chiếc máy photocopy đó sau 5 năm kể từ khi mua (kết quả làm tròn đến hàng phần chục).\\
	\shortans[oly]{$11,9$}
	\loigiai{
		Giá trị của máy photocopy sau $1$ năm sử dụng là
		$$T_{1}=50 \cdot 75 \%=37,5 \text{(triệu đồng).}$$
		Giá trị của máy photocopy sau $2$ năm sử dụng là
		$$T_{2}=T_{1} \cdot 75 \%=28,125 \text { (triệu đồng). }$$
		Giá trị của máy photocopy sau $3$ năm sử dụng là
		$$T_{3}=T_{2} \cdot 75 \%=21,0938 \text { (triệu đồng). }$$
		Giá trị của máy photocopy sau $4$ năm sử dụng là
		$$T_{4}=T_{3} \cdot 75 \%=15,8203 \text { (triệu đồng). }$$
		Giá trị của máy photocopy sau $5$ năm sử dụng là
		$$T_{5}=T_{4} \cdot 75 \%=11,8652 \text { (triệu đồng). }$$
		Chú ý: Tổng quát, giá trị của máy photocopy sau $n$ năm sử dụng là
		$$T_{n}=T_{1} \cdot(0,75)^{n-1} \text { (triệu đồng). }$$
	}
\end{ex}

\begin{ex}
	Một vi sinh đặc biệt $X$ có cách sinh sản vô tính kì lạ, sau một giờ thì đẻ một lần, đặc biệt sống được tới giờ thứ $n$ (với $n$ là số nguyên dương) thì ngay lập tức thời điểm đó nó đẻ một lần ra $2^n$ con $X$ khác, tuy nhiên do chu kì của con $X$ ngắn nên ngay sau khi đẻ xong lần thứ $2$, nó lập tức chết. Hỏi rằng, nếu tại thời điểm ban đầu có đúng $1$ con thì sau $5$ giờ có bao nhiêu con sinh vật $X$ đang sống?\\
	\shortans[oly]{$336$}
	\loigiai{
		Gọi $s_n$ là số lượng con $X$ được sinh ra tại thời gian thứ $n$, $n\in\mathbb{N}^*$ và $s_0=1$. Ta có
		\begin{align*}
			s_0 &=1,& s_1 &= 2 s_0= 2,\\
			s_2 &= 2s_1 + 4s_0 = 8, &s_3 &= 2s_2 + 4s_1 = 24,\\
			s_4 &= 2s_3 + 4s_2 = 80, &s_5 &= 2s_4 + 4s_3 = 256.
		\end{align*}
		Khi đó, sau $5$ giờ thì số con sinh vật $X$ đang sống bằng $s_4 + s_5 = 336$.
	}
\end{ex}

\begin{ex}
	Bác Hưng để $10$ triệu đồng trong tài khoản ngân hàng. Vào cuối mỗi năm, ngân hàng trả lãi $3 \%$ vào tài khoản của bác ấy, nhưng sau đó sẽ tính phí duy tri tài khoản hằng năm là $120$ nghìn đồng. Tìm số dư trong tài khoản của bác Hưng sau 4 năm (tính bằng triệu đồng và kết quả làm tròn đến hàng phần chục).\\
	\shortans[oly]{$10,8$}
	\loigiai{
		\begin{listEX}[1]
			\item  Vào cuối năm thứ nhất, số tiền trong tài khoản của bác Hưng là
			$$
			A_{1}=A_{0}(1+3 \%)-120000=1,03 A_{0}-120000 \text { (đồng). }
			$$
			Vào cuối năm thứ hai, số tiền trong tài khoản của bác Hưng là
			$$
			A_{2}=A_{1}(1+3 \%)-120000=1,03 A_{1}-120000 \text { (đồng). }
			$$
			Vào cuối năm thứ ba, số tiền trong tài khoản của bác Hưng là
			$$
			A_{3}=A_{2}(1+3 \%)-120000=1,03 A_{2}-120000 \text { (đồng). }
			$$
			Tương tự, vào cuối năm thứ $n(n \geq 1)$, số tiền trong tài khoản của bác Hưng là
			$$
			A_{n}=A_{n-1}(1+3 \%)-120000=1,03 A_{n-1}-120000 \text { (đồng). }
			$$
			\item Ta tính lần lượt $A_{1}, A_{2}, A_{3}, A_{4}$ :
			$$
			\begin{array}{ll}
				A_{1}=10180000 ; & A_{2}=10365400 ; \\
				A_{3}=10556362 ; & A_{4}=10753053 .
			\end{array}
			$$
			Như vậy, số dư trong tài khoản của bác Hưng sau $4$ năm là $10753053$ đồng.
		\end{listEX}
		
	}
\end{ex}

\Closesolutionfile{ans}


%\end{multicols}
\Closesolutionfile{ans}

